\section{Scope of the Problem}
This problem arises in essentially all JVM languages, all of which must interact with Java and have therefore chosen to also support method-local classes.

\begin{figure}
\begin{lstlisting}[numbers=left,language=c++,escapeinside={(*}{*)}]
class Object{};
template<typename X> class Foo : public Object {
    public: X x() { return x(); } };
template<typename X>
Foo<X>* m(Object* other, const X* x) {
    class MyFoo : public Foo<X> { public: X x() { return x; } };
    MyFoo* d = dynamic_cast<MyFoo*>(other);(*\label{line:cpp_err}*)
    if(d == nullptr) { return new MyFoo(); }
    return d;
}
\end{lstlisting}
\caption{C++ equivalent of Example~1.}
\label{fig:cpp}
\end{figure}

\subsection{Generic Reification} A natural suspicion is that this is ``simply'' a problem with generic erasure, that of course a language that cannot distinguish generic instantiations of a class at runtime will have difficulties of this sort.
There is some truth to this: Java is permitting a cast that it cannot actually check due to type erasure.
However there is more to it than that: Java will eagerly flag code such as \Q@(List<String>)(new LinkedList<Integer>())@.
So the language does intend to warn about uses of uncheckable casts.
\mjnote{\texttt{(List<String>)(new LinkedList<Integer>())} is an incompatible-types error, not an unchecked one. The contrast that matters here: an unchecked cast is a warning, an unchecked \texttt{instanceof} is an error.}
And while reified generics do provide enough information to check casts dynamically, only certain forms of reification
make correct casts straightforward.

In a language which implements parametric polymorphism by templating / heterogeneous compilation
(C++, C\#, or early experiments with generic Java~\cite{agesen1997adding}), the check is straightforward:
each instantiation of \Q@m@ with different generic parameters is compiled separately, resulting in separate
class declarations for each instantiation.
Consider reinterpreting Example~1 in C++ as in Figure \ref{fig:cpp}.
\texttt{gcc} reports for Line \ref{line:cpp_err}:
\begin{verbatim}
<source>:7:16: error: cannot 'dynamic_cast' 'other' (of type 'class Object*')
to type 'class m<Integer>(Object*, const Integer*)::MyFoo*' (source type is not polymorphic)
    7 |     MyFoo* d = dynamic_cast<MyFoo*>(other);
      |                ^~~~~~~~~~~~~~~~~~~~~~~~~~~
\end{verbatim}
\todo[inline]{not perfect, because this is C++ giving a static warning, not catching the error dynamically...}
Notice that the error message uses the instantiation of the method signature as a scope for a declaration of \Q@MyFoo@.
But heterogeneous compilation has subtle costs and wide-ranging consequences for language design, not just implementation,
and complex performance consequences~\cite{odersky2001two,kennedy2001design}. Even if this does fix the problem,
it is not universally appropriate.

An alternative, common conception of generic reification is to simply save the type parameters along with the
single implementations of classes, for it to be checked when needed. But note again that \Q@MyFoo@
has no generic parameters of its own. So a shallow check with this variant of reified generics is insufficient ---
each \Q@MyFoo@ would need to contain a representation of, or reference to, its parent type's generic parameters.
And because this phenomenon is not limited to a single step in the inheritance hierarchy, the cast must traverse and compare
a substantial amount of data --- not a mere comparison of virtual table pointers. This again has multiple implementation
approaches, neither ideal (chains of parent type instantiations hurt cache locality, inlining data can bloat small instances).

So there remains a need to articulate a solution to this problem where generics are implemented via erasure.
\mjnote{under erasure, Java's checked-cast rule is already enough once the dependency of \texttt{MyFoo} on \texttt{X} is visible in its type, and that is what funnelling provides. This paragraph can lead straight into Section~3.}



\todo[inline]{ah, we need to check Flix, where maybe we can cook up something roughly like this?}

//--------
MARCO: (6th time)
This is not about generic erasure
In a language with reified generics, if the language still decides that MyFoo is not generic,
no impact of those reified generics is going to be seen in the case (MyFoo)other
MyFoo extends Foo<X> is 'as generic' as
String extends Comparable<String>{}
A cast to String, would not look in the generics of Comparable<String>
A cast to MyFoo, would not look in the generics of Foo<X>
If subtyping has (nominal-)'reflexitity' then String <: String (no need to look at the generics in Comparable<String>, available or not.
A language with implicity 'type scope' that could conceptually add to the 'name' when types are identified in the type system could change this.
Eg: scala foo.Bar is a type and foo.Bar != beer.Bar //existential types on constants/values
Broken.<X>m.MyFoo  //existential types on generic-instantiated methods
//----------
COLIN: This is ambiguous: with reified generics, there are two implementation strategies, which yield different behaviours:
1. MyFoo allocated as a Foo<String> gets a different class table from MyFoo allocated as a Foo<Integer>. In this case, the cast will fail
2. All MyFoo will get the same inherited class table, but keep a separate field with the choice of X for Foo<X>, used for inherited methods. In this case, there are *two possible cast implementations*:
    a. Cast checks only vtable, and cast succeeds. Arguably a compiler that "treats MyFoo as non-generic" could lean towards this
    b. Cast checks both vtable and the runtime representation of the type parameter of Foo. Arguably a compiler that preserves this for inherited methods from Foo to work correctly would be likely to check this (since it would be checked in a cast to Foo<Integer>).

//----------
\mjnote{the shallow/deep question (2a/2b) only arises for a language that checks casts at run time instead of requiring the checked-cast premise; under modern Java it does not arise.}

This is a problem in other languages (cite cite) as shown in greater details in Rel Work.

Now, assume we have a readable small step reduction over a version of FJ rich enough to express the problem.
We will assume: 
- Simple generics
- A syntax like anonymous inner classes, but where we give them names. We can call them
named inner classes.
This imports the concept of method local inner classes within an expression based calculus.
//our model will also need casts + instanceof
\begin{lstlisting}
class A{ A onlyA()-> new A() }
class B{ B onlyB()-> new B() }
class Foo<X> extends Object{ X x() {return x();} }
public class Broken {
  <X> Foo<X> m(Object other, X x) ->
    if(other instanceof MyFoo)
      then ((MyFoo)other)
      else new MyFoo() extends Foo<X>{ X x(){ return x; } }//Bad
  }
main:
  new Broken().<B>m(new Broken().<A>m(new Object(),new A()),new B()).x().onlyB();
\end{lstlisting}


\begin{lstlisting}
class Foo<X> extends Object{ X x() {return x();} }
public class Broken {
  <X> Foo<X> m(Object other, X x) ->
    if(other instanceof MyFoo<X>) //unchecked instanceof (ill typed in our core)?
      then ((MyFoo<X>)other)//unchecked cast (ill typed in our core)?
      else new MyFoo<X>() extends Foo<X>{ public X x(){ return x; } }//good
  }
\end{lstlisting}
\mjnote{answered by Figure~\ref{fig:cast}: with \texttt{other : Object}, \texttt{other instanceof MyFoo<X>} is ill typed, because \texttt{X} is not determined by \texttt{Object} (as javac: \texttt{Object cannot be safely cast to MyFoo<X>}). With \texttt{other : Foo<X>} it is well typed and sound under erasure.}

Note how funnelling purpose is to insert the generics. Without funnelling, MyFoo is a non generic type, so the question about generic reification does not even arise.

small step reduction
\begin{verbatim}
01:  new Broken().<B>m(new Broken().<A>m(new Object(),new A()),new B()).x().onlyB();
02:  new Broken().<B>m(
    if(new Object() instanceof MyFoo<X>)
      then (MyFoo<A>) new Object()
      else new MyFoo<A>() extends Foo<A>{ public A x(){ return new A(); } }
      ,new B()).x().onlyB();
03:  new Broken().<B>m(
    if(false)
      then (MyFoo<A>) new Object()
      else new MyFoo<A>() extends Foo<A>{ public A x(){ return new A(); } }
      ,new B()).x().onlyB();
04:  new Broken().<B>m(
       new MyFoo<A>() extends Foo<A>{ public A x(){ return new A(); } }
      ,new B()).x().onlyB();
05:   (if(new MyFoo<A>() extends Foo<A>{ public A x(){ return new A(); } } instanceof MyFoo<B>)
       then (MyFoo<A>) new MyFoo<A>() extends Foo<A>{ public A x()->new A() }
       else new MyFoo<B>() extends Foo<B>{ public B x(){ return new B(); } }).x().onlyB()
06a:  (if(false)//reified gen
       then (MyFoo<B>) new MyFoo<A>() extends Foo<A>{ public A x(){ return new A(); } }
       else new MyFoo<B>() extends Foo<B>{ public B x(){ return new B(); } }).x().onlyB()
07a:  (new MyFoo<B>() extends Foo<B>{ public B x(){ return new B(); } }).x().onlyB()
08a:  new B().onlyB()
09a:  new B()//success and different semantic
06b:  (if(true)//Java
       then (MyFoo<B>) new MyFoo<A>() extends Foo<A>{ public A x(){ return new A(); } }
       else new MyFoo<B>() extends Foo<B>{ public B x(){ return new B(); } }).x().onlyB()
07b:  ((MyFoo<B>) new MyFoo<A>() extends Foo<A>{ public A x(){ return new A(); } }).x().onlyB()
08b:  (new MyFoo<A>() extends Foo<A>{ public A x()->new A() }).x().onlyB()
     //assuming cast and instanceof behave consistently (unchecked)
09b:  new A().x().onlyB()
10b:  new A().onlyB() //stuck!!
6c:  (if(err)//compile or at least run time
       then ((MyFoo<B>) new MyFoo<A>() extends Foo<A>{ public A x(){ return new A(); } }
       else new MyFoo<B>() extends Foo<B>{ public B x(){ return new B(); } }).x().onlyB()
\end{verbatim}

// Maybe above we want MyFoo:Object<A> and MyFoo:Object<B>?

\begin{lstlisting}
public class Broken {
  static <X> Object m(Object other, final X x, Function<X,String> f){
    class MyFoo{ public X x(){ return x; } }
    if(!(other instanceof MyFoo)){ return new MyFoo(); }
    MyFoo local= (MyFoo)other;//cast pass with no warnings
    System.out.println("Pass with " + x);
    System.out.println("old one was "+f.apply(local.x()));//throws ClassCastException
    return local;
  }
  public static void main(String[] a) {
    Object r1= Broken.<Integer>m("anything",5,i->""+i);
    Broken.<String>m(r1,"Hi",s->s);
  }
}
\end{lstlisting}

\begin{lstlisting}
public class Broken {
  static <X> Object m(Object other, final X x, Function<X,String> f){
    class MyFoo{ public X x(){ return x; } }
    if(!(other instanceof MyFoo)){ return new MyFoo(); }
    MyFoo local= (MyFoo)other;//cast pass with no warnings
    System.out.println("Pass with " + x);
    System.out.println("old one was "+f.apply(local.x()));//throws ClassCastException
    return local;
  }
  public static void main(String[] a) {
    Object r1= Broken.<Integer>m("anything",5,i->""+i);
    Broken.<String>m(r1,"Hi",s->s);
  }
}
\end{lstlisting}

\begin{lstlisting}
public class Broken {
  static <X> Object m(Object other, final X x, Function<X,String> f){
    class MyFoo{ public X x(){ return x; } }
    if(other instanceof MyFoo local){ //new instanceof+introduce var (pass with no warnings)
      System.out.println("Pass with " + x); //snekly `local` is MyFoo<Integer> even when this is Brokem.<String>m(..)
      System.out.println("old one was "+f.apply(local.x()));//throws ClassCastException    
     }
  return new MyFoo();//snekly it is MyFoo<Integer>
  }
  public static void main(String[] a) {
    Object r1= Broken.<Integer>m("anything",5,i->""+i);
    Broken.<String>m(r1,"Hi",s->s);
  }
}
\end{lstlisting}

// TODO: could talk about doing this with modern Java type-switch
\begin{lstlisting}[language=Java]
// Compiles without error, fails, on jdk 25.0.1
import java.util.function.Function;
class Broken {
  static <X> Object m(Object other, final X x, Function<X,String> f){
    class MyFoo{ public X x(){ return x; } }
    switch (other) {
        case MyFoo local: // implicit unchecked cast
            System.out.println("Pass with " + x); //snekly `local` is MyFoo<Integer> even when this is Brokem.<String>m(..)
            System.out.println("old one was "+f.apply(local.x()));//throws ClassCastException (on the return of f.apply)
            break;
        default:
            break;
    }
    if(other instanceof MyFoo local){ //new instanceof+introduce var (pass with no warnings)
      System.out.println("Pass with " + x); //snekly `local` is MyFoo<Integer> even when this is Brokem.<String>m(..)
      System.out.println("old one was "+f.apply(local.x()));//throws ClassCastException    
     }
  return new MyFoo();//snekly it is MyFoo<Integer>
  }
  public static void main(String[] a) {
    Object r1= Broken.<Integer>m("anything",5,i->""+i);
    Broken.<String>m(r1,"Hi",s->(s.length())+""); //Broken.<String>m(r1,"Hi",s->s);
  }
}
\end{lstlisting}

\begin{lstlisting}
//Example that breaks in a NEW WAY in the bytecode
// https://godbolt.org/z/sMnxPG1TE
import java.util.function.Function;
class Broken {
  static <X> Object m(Object other, final X x, Function<X,Integer> f){
    class MyFoo{ public X x(){ return x; } }
    switch (other) {
        case MyFoo local:
            System.out.println("Pass with " + x); //snekly `local` is MyFoo<Integer> even when this is Brokem.<String>m(..)
            System.out.println("old one was "+f.apply(local.x()));//throws ClassCastException 
            break;
        default:
            break;
    }
    if(other instanceof MyFoo local){ //new instanceof+introduce var (pass with no warnings)
      System.out.println("Pass with " + x); //snekly `local` is MyFoo<Integer> even when this is Brokem.<String>m(..)
      System.out.println("old one was "+f.apply(local.x()));//throws ClassCastException    
     }
  return new MyFoo();//snekly it is MyFoo<Integer>
  }
  public static void main(String[] a) {
    Object r1= Broken.<Integer>m("anything",5,i->i);
    Broken.<String>m(r1,"Hi",s->s.length());
  }
}
// Above seems to fail but without explicit checkcast, failing one of the virtual method calls
\end{lstlisting}
With investiagion in bytecode and trying both s->s.length() and String::length we think that in order to save on
bytecode size, our JVM enviroment under test adds casts in all lambda calls so that meth refecences lambda calls do not need specialized bytecode (eg MyInterface a=String::length)

% interface Shape{ Square toSquare(); }
% class Square extends Shape{Square toSquare(){ return this; } }
%
%
%
% public class Broken {
%   static <X> Foo<X,?> m(Top<?> other, final X x){
%     class MyFoo extends Foo<X,MyFoo>{
%       public X x(){ return x; }
%       public MyFoo myFooOrElse(MyFoo myFoo){ return this; }
%     }
%     return other.selfOrElse(new MyFoo());
%   }
%   public static void main(String[] a) {
%     Foo<Integer> r1= Broken.<Integer>m(new Top<String>(),5);
%     Foo<String> r2= Broken.<String>m(r1,"Hi");
%     System.out.println(r1.x());//5
%     System.out.println(r2.x());//Throws ClassCastException: ..
%   }
% }
% class Top<S>{
%   public S selfOrElse(S s){ return s; }
%   }
% class Foo<X,S> extends Top<S>{ public X x() {return x();} }
%
%
